\chapter{Conclusion}

This work aimed to explore the intrinsic connection between the convolutional structure of neural
networks (CNNs) and the emergence of preferred directions of vulnerability in the
Fourier domain. More specifically, our approach sought to determine the extent to which the
mathematical properties inherent in convolutional layers, beyond biases associated with the data or
the training process, can explain the existence of universal adversarial perturbations.

To this end, after establishing the general framework of CNNs and adversarial attacks, we
reconstructed and analyzed the theoretical reasoning proposed by Tsuzuku and
Sato~\cite{tsuzuku2019structural}, initially using a simplified linear model.

To extend these theoretical results to actual CNN architectures, where
ReLU activations and \textit{max-pooling} operations prevent any
global diagonalization, we introduced local linearization through
analysis of the model's Jacobian. This theoretical extension proved crucial
in justifying the application of results based on linearity to
fundamentally nonlinear architectures.

Together, these observations clarify the mechanism by which a
perturbation based on a single Fourier frequency can affect a convolutional
network so severely. The successive amplification of this perturbation through
the different layers provides a compelling structural explanation for the existence of
universal perturbations.

Our experimental approach supports this intuition. The \textit{heatmaps}
obtained for \texttt{VGGSmallBN} on the CIFAR-10 and
CIFAR-100 datasets reveal highly heterogeneous frequency sensitivity.
Certain frequencies produce significantly higher \textit{fool ratios},
showing that perturbations of equal amplitude do not have an
equivalent impact: their direction in input space is a determining factor.
Comparison with random noise of the same amplitude
fully corroborates this asymmetry.

The results also highlight the expected influence of classification task
complexity on model robustness. Overall, the network proves more
vulnerable on CIFAR-100 than on CIFAR-10. This vulnerability, examined in greater detail
in Appendix~\ref{annexe:marges-logits}, is consistent with the increased difficulty of the problem:
with more classes, decision boundaries become more numerous and more closely
spaced because of the visual similarity between categories, allowing a perturbation
to change the prediction more easily. Nevertheless, despite this increased vulnerability,
sensitivity retains a structured topology in which certain Fourier directions remain dominant.

The scope of these results should nevertheless be interpreted cautiously. Our experimental reproduction
is limited to a VGG-type architecture and low-resolution images. Moreover, the
most sensitive frequencies were selected from \textit{heatmaps} computed on
the test set. We therefore cannot claim that all convolutional networks exhibit
exactly the same vulnerability structure. Within the setting studied here, however, Fourier
directions appear to be a useful tool for better understanding certain sensitivities of CNNs.

This analysis can also be related to model explainability. When a
particular frequency substantially changes the network's predictions, this indicates that this component
of the image plays an important role in its decision. In other words, sensitive frequencies
are useful not only for constructing adversarial perturbations: they also provide clues
about the type of information the network uses to classify images.

A possible continuation of this work would therefore be to use these sensitive frequencies both to strengthen
models and to interpret them more effectively. For example, they could guide new
forms of data augmentation, limit certain internal amplifications, or support comparisons of
more robust architectures, particularly those incorporating anti-aliasing mechanisms.

This direction would connect three important aspects: adversarial robustness,
resistance to variations in the data, and interpretation of the network's decisions.
The goal would therefore be not only to show that certain Fourier frequencies are dangerous
for the model, but also to understand how identifying them can help build networks
that are more stable, more robust, and more interpretable.

